% introduction_critiques.tex
\subsection{Critiques and Boundary Conditions}

The same capabilities, however, threaten classical psychometrics' response data: the autonomous online ``synthetic respondent'' of \citet{westwood-2025}, a coherent demographic persona with memory of prior answers, passed 99.8\% of 6{,}000 attention-check trials, gave largely internally consistent (correctly reverse-scored) responses to established scales, strategically declined ``reverse shibboleth'' probes of superhuman capability, shifted aggregate polling on instruction, and---most insidiously---inferred researchers' hypotheses from the materials and biased answers toward them by 22.2 percentage points over the original in a replicated classic survey experiment. Westwood's conclusion---a coherent response can no longer be assumed human---cuts two ways: synthetic coherence is cheap and not validity, no substitute for human evidence; yet semantic factor analysis treats embeddings as measures of item meaning, not simulated respondents, and semantic--behavioral convergence as the empirical question \texttt{semanticfa} operationalizes.

Analyzing item language also faces the sharpest critique of language-based assessment: for \citet{uher-2025}, psychometrics is pragmatic quantification, not measurement---discrimination and prediction with no traceable connection to the quantities to be measured---its cardinal error mistaking judgments of verbal statements for measurements of the phenomena they describe. Standardized items have no unitary meanings in use---each rater construes a local sample of an item's collective field of meaning---so rating-based nomological networks may substantially reflect the items' \emph{inbuilt semantics}, the relations among their words, rather than the phenomena, and language-model scale design or refinement risks perpetuating the error. The package does not contest the critique: its declared object is the inbuilt semantics---the semantic proposal encoded in the item set---and, with no respondent, it measures no attribute of any person and asserts no attribute-to-score causal path in the sense of \citet{borsboom-etal-2004}; it can characterize the proposal and compare it with other evidence, never validate a construct or instrument, but it makes the critique's central conjecture estimable scale by scale: the coherence statistics below quantify how much response structure the wording alone anticipates. High coherence vindicates the critique measurably and grounds pre-data structural screening; low coherence marks behavioral covariance semantics cannot explain---a candidate signature of the divergence conditions above or the language--phenomena breaks \citet{uher-2025} describes. Either way, what items say versus how people answer becomes a reported instrument property, not an unexamined confound; Uher's own constructive program---language technologies on unrestricted verbal self-descriptions---complements rather than competes with this use.
